% ============================================================
% REVISED §IV-A: Data Source and Generation Pipeline
% Replaces lines 384–441 of main.tex
% ============================================================

\subsection{Data Source and Generation Pipeline}
\label{sec:data}

\subsubsection{Simulation environment and flight characteristics}

We generate the experimental dataset using the MATLAB UAV Toolbox
(\textit{UAV Toolbox + Navigation Toolbox + Sensor Fusion and Tracking Toolbox})
because it provides physics-grounded sensor models for GPS, IMU, and INS fusion
that are absent in software-in-the-loop flight stacks.
Four mission profiles modelling power-line inspection tasks are generated:
power-line scan, rectangular area survey, tower orbit, and long-distance transit.
For each flight we simulate a complete waypoint mission executed by a
\texttt{uavWaypointFollower} controller at a cruise speed of 5 m/s
and a flight altitude of 50 m AGL.

The sensor pipeline combines:
\begin{itemize}
  \item \textbf{GPS model:} \texttt{gpsSensor} with horizontal position noise
        $\sigma_{xy}=\SI{1.5}{\metre}$, vertical noise $\sigma_z=\SI{3.0}{\metre}$,
        and velocity noise $\sigma_v=\SI{0.05}{\metre\per\second}$ (1-$\sigma$);
  \item \textbf{IMU model:} \texttt{imuSensor} with accelerometer noise
        $\sigma_a=\SI{0.02}{\metre\per\second\squared}$ and gyroscope noise
        $\sigma_\omega=\SI{0.001}{\radian\per\second}$;
  \item \textbf{INS fusion:} \texttt{insfilterAsync} for tight GPS/IMU fusion,
        producing a smoothed 9-dimensional navigation state vector at 10 Hz.
\end{itemize}

The 9-d state vector consists of ENU position $(x,y,z)$, velocity $(v_x,v_y,v_z)$,
and linear acceleration $(a_x,a_y,a_z)$. This matches the feature set used in
our earlier formulation (Section~\ref{sec:framework}) and is sufficient to
characterise kinematic disturbances from GNSS spoofing.

\textbf{Simulation fidelity and limitations.}
The MATLAB sensor models use additive Gaussian white noise and do not include
multipath propagation, ionospheric delay, or receiver clock dynamics.
We assess the resulting \emph{sim-to-real gap} quantitatively in the
cross-domain evaluation (Section~\ref{sec:cross_domain}).

\subsubsection{Dataset scale and split}

We generate \num{3000} independent flights totalling approximately
$3000 \times 60\,\text{s} = 180{,}000$ seconds of simulated flight time
(Table~\ref{tab:dataset}).
Flights are partitioned at the \emph{flight level} to prevent overlapping windows
from appearing in more than one split; within each split, 40\% of flights are
randomly selected for attack injection.

\begin{table}[!t]
\caption{MATLAB-generated dataset statistics.}
\label{tab:dataset}
\centering
\begin{tabular}{lccc}
\hline
\textbf{Split} & \textbf{Flights} & \textbf{Attacked} & \textbf{Windows} \\
\hline
Train      & 1800 & 739 (41\%) & 212{,}758 \\
Validation & 600  & 237 (40\%) & 69{,}993  \\
Test       & 600  & 229 (38\%) & 69{,}636  \\
\hline
\textbf{Total} & \textbf{3000} & \textbf{1205 (40\%)} & \textbf{352{,}387} \\
\hline
\end{tabular}
\end{table}

The test set flights are drawn from mission seeds, origin coordinates, and attack
parameter ranges that are held out from training, providing a mild out-of-distribution
evaluation within the simulated domain.

\subsubsection{Train-only data augmentation}

If needed to address class imbalance, TimeGAN-based augmentation is applied
\emph{exclusively to the training split}. The validation and test sets contain
only original simulated flights and are never subjected to any synthetic augmentation.
This constraint is enforced programmatically; an assertion in the data pipeline
raises an error if any augmented flight ID appears in the evaluation splits.
The data-flow is summarised in Fig.~\ref{fig:dataflow}.

\begin{figure}[!t]
\centering
\includegraphics[width=0.9\columnwidth]{dataflow.pdf}
\caption{Data-flow diagram.  Augmentation is applied only to the training split;
evaluation splits are strictly isolated from synthetic data.}
\label{fig:dataflow}
\end{figure}
